% ==========================================
\section{Sistema Attuale}
% ==========================================

% Oggi uno studente che vuole presentarsi a un'audizione, candidarsi per una borsa di studio
% o collaborare con un ente esterno si trova a raccogliere manualmente file sparsi,
% email con allegati, portfolio in PDF e curriculum aggiornati a mano. 
% Ogni contesto richiede una presentazione diversa, ma ricostruirla da zero ogni volta è 
% dispersivo e poco efficace.

L'analisi dello scenario tecnologico corrente evidenzia l'assoluta mancanza di una piattaforma integrata e istituzionale dedita alla valorizzazione dei flussi di lavoro dei candidati AFAM. La situazione attuale \`e caratterizzata da un modello fortemente decentralizzato e destrutturato, all'interno del quale gli studenti sono costretti a gestire la propria presentazione professionale in modo manuale e frammentato.

Attualmente, la sottomissione di un percorso artistico per audizioni, borse di studio o collaborazioni esterne richiede l'utilizzo concorrente di canali eterogenei: la posta elettronica per le comunicazioni formali, i servizi di cloud storage commerciali per l'invio di file multimediali ad alta risoluzione (registrazioni audio, tracce video, spartiti digitalizzati) e la compilazione estemporanea di documenti di testo per i curricula. Questo approccio comporta un elevato overhead di gestione, poich\'e l'allineamento e l'aggiornamento dei file rimangono a totale carico dello studente, incrementando il rischio di inviare documentazione obsoleta o non coerente.

Sotto il profilo del controllo e della sicurezza, il sistema attuale soffre di una totale asimmetria informativa. Una volta spedito il materiale via email o tramite link di terze parti, l'utente perde qualsiasi tracciamento sul ciclo di vita del file: non \`e possibile verificare se i valutatori abbiano effettivamente visionato le opere, n\'e revocare l'accesso ai documenti in caso di variazioni contrattuali o di ripensamenti sulla privacy. Inoltre, la tutela della propriet\`a intellettuale dei materiali originali risulta gravemente compromessa dalla mancanza di un ambiente di fruizione controllato dall'istituzione.